\documentclass[11pt,a4paper]{article}

\usepackage[T1]{fontenc}
\usepackage[utf8]{inputenc}
\usepackage{lmodern}
\usepackage[french]{babel}
\usepackage[a4paper,margin=2.3cm]{geometry}
\usepackage{microtype}
\usepackage{xcolor}
\usepackage{enumitem}
\usepackage{titlesec}
\usepackage{hyperref}
\usepackage[most]{tcolorbox}

% ---- palette ----
\definecolor{ink}{HTML}{1c2330}
\definecolor{accent}{HTML}{2f6e4c}
\definecolor{accent2}{HTML}{8a5a1c}
\definecolor{soft}{HTML}{eef3ef}
\definecolor{softline}{HTML}{cfe0d4}
\definecolor{warn}{HTML}{7a2418}
\definecolor{warnsoft}{HTML}{fbeeea}
\definecolor{warnline}{HTML}{e7c3ba}

\hypersetup{colorlinks=true, linkcolor=accent, urlcolor=accent2, pdftitle={Étude — Progression long-terme (tactique TB \& RPG)}}

\titleformat{\section}{\Large\bfseries\color{ink}}{\thesection.}{0.5em}{}
\titleformat{\subsection}{\large\bfseries\color{accent}}{\thesubsection}{0.6em}{}
\titlespacing{\section}{0pt}{1.4em}{0.5em}
\setlist[itemize]{leftmargin=1.2em,itemsep=2pt,topsep=2pt}

% ---- callouts ----
\newtcolorbox{principe}{enhanced,breakable,colback=soft,colframe=accent,
  boxrule=0.4pt,left=8pt,right=8pt,top=5pt,bottom=5pt,arc=2pt,
  fonttitle=\bfseries,coltitle=accent,title=Principe}
\newtcolorbox{antipat}{enhanced,breakable,colback=warnsoft,colframe=warn,
  boxrule=0.4pt,left=8pt,right=8pt,top=5pt,bottom=5pt,arc=2pt,
  fonttitle=\bfseries,coltitle=warn,title=Anti-pattern}
\newtcolorbox{optionbox}[1]{enhanced,breakable,colback=white,colframe=softline,
  boxrule=0.6pt,left=8pt,right=8pt,top=5pt,bottom=5pt,arc=2pt,
  fonttitle=\bfseries,coltitle=ink,title={#1}}

\newcommand{\dv}{\,\textsuperscript{\textcolor{accent2}{\textdagger}}}
\newcommand{\jeu}[1]{\textbf{#1}}
\newcommand{\compromis}[1]{\par\smallskip\textit{\textcolor{warn}{Compromis :}} #1}

\title{\vspace{-1.5em}\textbf{\color{ink}\Huge Rendre la progression motivante}\\[0.3em]
\large\color{accent}Étude de design --- tactique tour-par-tour \& RPG\\[0.2em]
\normalsize\color{ink}Éviter le treadmill XP/loot et la linéarité, sans gonfler le monde}
\author{}
\date{}

\begin{document}
\maketitle
\vspace{-3em}

\begin{center}
\small\itshape Panorama de $\sim$25 jeux, organisé en \textbf{menus d'options} par couche (macro / micro / tension).\\
Darkest Dungeon et XCOM n'y figurent que comme deux exemples parmi beaucoup d'autres.
\end{center}

\begin{tcolorbox}[colback=soft,colframe=softline,boxrule=0.5pt,arc=2pt,left=8pt,right=8pt,top=5pt,bottom=5pt]
\small\textbf{Niveaux de confiance.} La majorité des mécaniques décrites sont confirmées par des wikis/guides communautaires et la connaissance établie de ces jeux. Les points marqués \dv{} sont corroborés par une \textbf{source de développeur} (post-mortem, talk GDC, interview). Les \textbf{compromis} et la synthèse transversale relèvent de l'\textbf{analyse de design}, pas de faits sourcés. Sources en annexe.
\end{tcolorbox}

\tableofcontents

% ======================================================================
\section{Diagnostic : pourquoi la boucle « XP / loot / voyage » lasse}

Le treadmill classique empile \emph{des nombres} : +1 stat par niveau, du loot incrémental, des ennemis qui montent au même rythme que vous. Trois défauts structurels :

\begin{itemize}
\item \textbf{L'inflation verticale dévalorise le passé.} Multiplier les gains d'XP dévalue l'effort accumulé, exactement comme une inflation monétaire : le progrès devient une illusion court-terme. C'est l'écueil n\textsuperscript{o}1.
\item \textbf{Pas de décision = pas d'intérêt.} Si « monter » se résume à répéter une boucle pour franchir un seuil, on \emph{exécute} ; on ne \emph{décide} pas.
\item \textbf{Pas d'enjeu = pas d'attachement.} Si rien d'irremplaçable ne peut être perdu, la progression n'a aucun poids.
\end{itemize}

\begin{principe}
Les trois vrais moteurs de motivation long-terme, du plus fort au plus faible : \textbf{(1) l'enjeu / l'attrition} (on peut perdre), \textbf{(2) la décision} (chaque palier ferme des portes), \textbf{(3) l'accumulation de stats} (le plus faible — à n'utiliser que comme \emph{support} des deux autres). La suite décline un \emph{menu d'options} pour chacun.
\end{principe}

% ======================================================================
\section{Couche MACRO --- structurer la campagne et le monde}
\emph{Objectif : créer la non-linéarité par le \textbf{choix}, pas par la \textbf{taille} ; et empêcher le farm par la \textbf{pression}.}

\begin{optionbox}{Option A1 --- L'horloge / la menace montante}
Une pression qui s'aggrave avec le temps rend le farm contre-productif : attendre coûte.
\begin{itemize}
\item \jeu{XCOM 2} : le \emph{Projet Avatar} est un compte à rebours mondial ; plein, déclenche la défaite. \jeu{WOTC} ajoute les \emph{Élus} qui montent en puissance entre les rencontres et la \emph{fatigue} des soldats (on ne peut pas aligner sa A-team indéfiniment).
\item \jeu{Phoenix Point} : la \emph{Brume} Pandoréenne s'étend en temps réel ; les nids non traités deviennent des Citadelles plus dangereuses.
\item \jeu{Invisible, Inc.} : l'\emph{Alarme} monte à chaque tour (modelée sur la faim des roguelikes\dv) + compte à rebours global de 72 h. Camper est mortel.
\item \jeu{FTL} : la flotte rebelle vous poursuit balise par balise ; s'attarder pour farmer, c'est se faire rattraper.
\item \jeu{Battle Brothers} : des \emph{crises} montantes obligatoires (invasions) reconfigurent la carte et donnent une fin claire.
\end{itemize}
\compromis{trop de pression = anxiété ; il faut des respirations. Et des joueurs experts apprennent à \emph{retarder} l'horloge optimalement (XCOM 2), réduisant la tension prévue.}
\end{optionbox}

\begin{optionbox}{Option A2 --- La carte à nœuds (curation contre dilution)}
Toute la non-linéarité, zéro traversée vide : un graphe de choix de route.
\begin{itemize}
\item \jeu{Slay the Spire} : actes de $\sim$15 paliers, on choisit un nœud parmi 1--3 au-dessus, sans retour (combat / élite / ? / repos / marchand / boss).
\item \jeu{FTL} : balises par secteur ; on ne peut pas tout visiter, on lit le risque/récompense de chaque saut.
\item \jeu{Into the Breach} : 4 missions choisies sur $\sim$7 par île, ordre libre ; 2 îles suffisent pour la finale.
\item \jeu{Darkest Dungeon} : boucle hub (le Hameau persistant) $\leftrightarrow$ expéditions hebdomadaires dans 4 régions.
\item \jeu{The Banner Saga} : road-trip de caravane en chapitres, ponctué de hubs de repos.
\end{itemize}
\compromis{moins d'immersion « monde » et d'exploration libre qu'un DOS2 ; forte dépendance à la RNG du draft/des routes.}
\end{optionbox}

\begin{optionbox}{Option A3 --- Le geoscape : coût d'opportunité}
La « carte » est un espace de décisions simultanées, pas un terrain à parcourir.
\begin{itemize}
\item \jeu{XCOM} : 3 à 7 missions/projets concurrents en permanence ; on en choisit un, les autres font monter la panique / un Dark Event. On ne peut pas tout faire.
\item \jeu{Phoenix Point} : des dizaines de havres de 3 factions lancent des requêtes simultanées ; portée limitée des avions $\Rightarrow$ triage. Aider une faction en fâche une autre.
\item \jeu{Jagged Alliance 2} : contrôle de secteurs, milices, mines qui financent la guerre.
\end{itemize}
\end{optionbox}

\begin{optionbox}{Option A4 --- La libération de territoire comme colonne vertébrale}
\begin{itemize}
\item \jeu{Jagged Alliance 2/3} : capturer des secteurs \emph{est} la progression ; les milices auto-défendent ce qu'on a pris (on s'étend sans escorter chaque case), mais l'ennemi contre-attaque.
\item \jeu{Mount \& Blade II : Bannerlord} : fiefs, ateliers, caravanes financent une armée dont la \emph{composition} (haut tier) prime sur le nombre.
\end{itemize}
\compromis{sans crise scénarisée ni fin imposée (Bannerlord), la fin de partie « à plat » une fois qu'on surclasse tout le monde.}
\end{optionbox}

\begin{optionbox}{Option A5 --- Régions cloisonnées à sens unique}
Anti-bloat : un grand monde découpé en actes irréversibles, denses et réactifs.
\begin{itemize}
\item \jeu{Divinity: Original Sin 2} : 5 actes / régions ; passer un acte est définitif (contenu manquable sauf stocké au camp). Évite la map unique diluée.
\item \jeu{Wasteland 2/3} : réactivité « macro », quêtes concentrées par lieu, scope resserré sur une région (le Colorado pour W3).
\end{itemize}
\compromis{beaucoup de contenu manquable $\Rightarrow$ forte incitation aux guides en première partie.}
\end{optionbox}

\begin{optionbox}{Option A6 --- Procédural borné + « golden path » fait-main}
\jeu{XCOM} (cartes procédurales + missions d'histoire obligatoires), \jeu{Darkest Dungeon} (donjons procéduraux, bosses fait-main), \jeu{Into the Breach} (îles modulaires) : on combine variété/rejouabilité \emph{et} beats narratifs maîtrisés.
\end{optionbox}

\begin{optionbox}{Option A7 --- Le scaling auto contre la dilution}
\begin{itemize}
\item \jeu{Wartales} : le \emph{nombre} d'ennemis scale avec votre nombre de combattants $\Rightarrow$ « zerguer » des corps faibles invoque plus d'ennemis : croissance auto-limitée, on pousse vers une escouade lean.
\item \jeu{Bannerlord} : \emph{aucun} scaling — la puissance vient de la composition (petite armée haut-tier $>$ horde de bleus).
\end{itemize}
\compromis{le level-scaling \emph{généralisé} (cf. anti-pattern) annule la sensation de progrès ; le scaling « par le nombre » de Wartales l'évite en restant tactique.}
\end{optionbox}

\begin{optionbox}{Option A8 --- L'absence \emph{volontaire} de pression}
\jeu{Tactical Breach Wizards} : pas de limite de tours, système d'\emph{undo}, pas de geoscape — chaque mission est un puzzle auto-contenu, la pression est \emph{opt-in} (objectifs bonus). Maximise l'accessibilité et la combo ; sacrifie la tension de survie. Une option assumée, à l'opposé de tout le reste.
\end{optionbox}

\begin{antipat}
\textbf{Le monde ouvert gonflé.} Grande carte + points d'interrogation + listes de tâches + longs trajets = fatigue d'exploration. \emph{La taille n'est pas du contenu} : ce n'est pas l'open world qui lasse mais le monde \emph{vide}. La densité et la réactivité battent l'échelle.
\end{antipat}

% ======================================================================
\section{Couche MICRO --- montée en puissance, builds, économie}
\emph{Objectif : remplacer « +1 stat » par des \textbf{décisions} et des \textbf{paliers} qui changent \emph{comment} on joue.}

\begin{optionbox}{Option B1 --- Choix exclusifs verrouillés (au lieu de croissance plate)}
Chaque palier = une \emph{bifurcation}, pas un incrément.
\begin{itemize}
\item \jeu{Slay the Spire} : on choisit 1 carte sur 3 (ou aucune) ; le deck \emph{est} le build, construit par engagements successifs.
\item \jeu{XCOM} : à chaque grade, une capacité « A ou B » (jamais les deux) $\Rightarrow$ deux soldats de même classe divergent.
\item \jeu{Wartales} : 7 paliers par classe, un skill choisi par palier verrouille les autres.
\item \jeu{The Banner Saga} : promotions \emph{irréversibles} (gâtées par les kills, pas l'XP).
\item \jeu{Fire Emblem} : promotions/classes ; \jeu{Three Houses} : choix de maison/cours durables.
\end{itemize}
\end{optionbox}

\begin{optionbox}{Option B2 --- Pas d'XP du tout : la puissance \emph{structurelle}}
\jeu{Into the Breach} (puissance = mechs/armes qu'on \emph{débloque et choisit}), \jeu{FTL} (décisions sur le scrap de \emph{cette} run), \jeu{Slay the Spire} (le deck \& les reliques). Leçon transverse : \textbf{l'XP n'est qu'une option parmi d'autres pour faire monter}, souvent la plus paresseuse.
\end{optionbox}

\begin{optionbox}{Option B3 --- Learn-by-doing vs niveaux discrets}
\begin{itemize}
\item \emph{À l'usage} : \jeu{Jagged Alliance 2} (le tir monte le tir), \jeu{Bannerlord} (skills par usage, plafonnés par Focus/Attributs). Immersif mais opaque, invite au grind.
\item \emph{Discret} : \jeu{Jagged Alliance 3} (niveaux 1--10, perks débloqués à des seuils), \jeu{XCOM} (grades). Lisible, plafonné, déplace la puissance vers les \emph{choix} de perks/équipement.
\end{itemize}
\end{optionbox}

\begin{optionbox}{Option B4 --- Le potentiel fixé au recrutement (rareté humaine)}
\jeu{Battle Brothers} : à l'embauche, des « talents » aléatoires \emph{élargissent} les fourchettes de montée d'un combattant $\Rightarrow$ le plafond est décidé au recrutement, pas farmable. L'équipement compte autant que les stats.
\end{optionbox}

\begin{optionbox}{Option B5 --- Le build par composition d'équipe}
\jeu{Wasteland 2/3} (peu de points $\Rightarrow$ 7 Rangers spécialisés, complémentaires), \jeu{Divinity: OS2} (classes = simples templates, tout respéçable ; les sorts naissent de la \emph{combinaison} d'écoles, ex. Hydro+Necro). La profondeur est dans l'\emph{assemblage}, pas le niveau.
\end{optionbox}

\begin{optionbox}{Option B6 --- Multiclassage / archetypes / voies transformatives}
\jeu{Pathfinder: WOTR} : 10 \emph{Mythic Paths} (un run Ange et un run Filou « n'ont presque rien en commun » à l'acte 5) ; certaines voies \emph{fusionnent} magie et mythique. Builds transformatifs, pas juste « plus gros ».
\end{optionbox}

\begin{optionbox}{Option B7 --- La monnaie unique qui crée un dilemme permanent}
\begin{itemize}
\item \jeu{The Banner Saga} : le \emph{Renown} sert à la fois à monter de niveau, acheter des objets \emph{et} nourrir la caravane $\Rightarrow$ « renforcer les héros vs survivre » à chaque instant.
\item \jeu{Shadowrun: Dragonfall} : le \emph{Karma} paie la montée \emph{et} les Étiquettes (dialogues/contournements) ; chasser les étiquettes coûte au rôle de combat $\Rightarrow$ vrais coûts d'opportunité.
\end{itemize}
\end{optionbox}

\begin{optionbox}{Option B8 --- Loot rare, situationnel, à double tranchant}
\jeu{Battle Brothers} (objets nommés rares, dans des camps lointains), \jeu{Darkest Dungeon} (trinkets puissants \emph{avec un malus}) : équiper devient une décision, pas un empilement. Économie tendue (\jeu{JA2}, \jeu{Wartales}) : chaque achat est un arbitrage.
\end{optionbox}

\begin{optionbox}{Option B9 --- La customisation comme progression d'identité}
\jeu{XCOM}\dv{} : surnoms par classe, bios, apparence $\Rightarrow$ l'esprit du joueur « comble le récit », l'unité jetable devient un personnage. La customisation \emph{est} une forme de progression, distincte du stat-bloat.
\end{optionbox}

\begin{antipat}
\textbf{Le +1 stat linéaire + le level-scaling des ennemis.} Sans bifurcation ni nouveau « verbe », aucune identité ni décision ; et si tout monte avec vous, vous ne vous sentez \emph{jamais} plus fort (le progrès s'annule). Idem pour le \emph{robinet à loot} (slot-machine) qui supprime toute décision par abondance.
\end{antipat}

\begin{principe}
Privilégier l'\textbf{horizontal} (nouveaux outils/verbes qui ouvrent des tactiques) au \textbf{vertical} (nombres plus gros). « Monter », c'est \emph{se spécialiser} (fermer des options) plutôt qu'« avoir plus » : la spécialisation crée de l'identité et donne envie de rejouer.
\end{principe}

% ======================================================================
\section{Couche TENSION / ENJEUX --- ce qui donne \emph{envie} de continuer}
\emph{Le moteur le plus puissant, et le plus négligé.}

\begin{optionbox}{Option C1 --- La mort permanente, en gradients}
\begin{itemize}
\item \emph{Stricte} : \jeu{Fire Emblem} (mode classique), \jeu{Battle Brothers}, \jeu{Darkest Dungeon}.
\item \emph{Downed-then-dead} : \jeu{Jagged Alliance 3} (on saigne, un médecin peut stabiliser), \jeu{Wartales} (perte si l'unité \emph{fuit de terreur}).
\item \emph{Asymétrique} : \jeu{Bannerlord} (les troupes meurent, les héros sont seulement K.O.).
\item \emph{Optionnelle} : \jeu{FE: Three Houses} (Classic/Casual + rembobinage), ce qui \emph{dilue} l'attachement que la permadeath stricte produisait.
\end{itemize}
\end{optionbox}

\begin{optionbox}{Option C2 --- Blessures et séquelles permanentes}
\jeu{Battle Brothers} (œil perdu, lésion cérébrale après un coup encaissé), \jeu{Wildermyth}\dv{} (mutilations remplacées par des transformations, vieillissement, mort de vieillesse) : les marques s'inscrivent \emph{durablement} et deviennent l'histoire du personnage.
\end{optionbox}

\begin{optionbox}{Option C3 --- L'attrition (psychologique, morale, financière)}
\begin{itemize}
\item \emph{Psychologique} : \jeu{Darkest Dungeon}\dv{} — le \emph{Stress} est conçu \emph{explicitement} pour la tension émotionnelle, pas le min-maxing ; ressource persistante à deux faces (à 100 : Affliction handicapante \emph{ou}, plus rarement, Vertu héroïque). La récupération passe par des activités de ville \emph{en nombre limité} $\Rightarrow$ impossible de farmer le repos, on \emph{fait tourner le roster}. Un exemple fort, mais \textbf{un parmi d'autres}.
\item \emph{Morale / désertion} : \jeu{Wartales}, \jeu{Battle Brothers} (humeur basse $\Rightarrow$ désertions), \jeu{Bannerlord}.
\item \emph{Financière} : presque tous les jeux de mercenaires utilisent la \textbf{paie qui scale avec la taille/qualité} comme frein principal au snowball (\jeu{JA2/3}, \jeu{Battle Brothers}, \jeu{Wartales}, \jeu{Bannerlord}).
\end{itemize}
\end{optionbox}

\begin{optionbox}{Option C4 --- L'échec qui fait avancer (« failing forward »)}
\jeu{XCOM} (rater une mission fait monter l'Avatar mais la partie continue, transformée), \jeu{Battle Brothers} / \jeu{Darkest Dungeon} (un wipe est un revers de campagne, pas un load). Quand l'échec produit du \emph{contenu}, la tension devient jouable au lieu d'être contournée par le save-scum.
\end{optionbox}

\begin{optionbox}{Option C5 --- Les choix irréversibles et exclusifs (la vraie non-linéarité)}
\begin{itemize}
\item \jeu{Divinity: OS2} : Origins aux objectifs exclusifs ; un seul devient Divin (en coop, la fin se joue en trahison).
\item \jeu{Pathfinder: WOTR/Kingmaker} : voies mythiques exclusives ; états d'échec du royaume (game over si une stat tombe).
\item \jeu{Wasteland 2} : Highpool \emph{vs} Ag Center appellent à l'aide en même temps — l'un est détruit pour toute la partie. \jeu{Wasteland 3} : allégeance de faction exclusive.
\item \jeu{Shadowrun}, \jeu{The Banner Saga} : décisions qui ferment des portes, importées de chapitre en chapitre.
\end{itemize}
\compromis{« non-linéaire » $\neq$ « grand » : c'est \emph{des choix qui comptent et qu'on ne peut pas annuler}. Mais : beaucoup de contenu jamais vu en une partie, et risque d'« illusion de choix » si les conséquences sont cosmétiques.}
\end{optionbox}

\begin{optionbox}{Option C6 --- La narration émergente}
Des \emph{systèmes qui fabriquent des histoires} que le joueur raconte.
\begin{itemize}
\item \jeu{Crusader Kings 2/3}\dv{} : « des histoires sans intrigue ni script » (Fåhraeus, GDC) ; les \emph{traits} des agents autonomes (opinions $-100$ à $+100$) génèrent intrigues, trahisons, successions imprévues.
\item \jeu{XCOM}\dv{} / \jeu{Darkest Dungeon}\dv{} : permadeath + identité $\Rightarrow$ récit personnel ; DD suit l'historique d'affliction de chaque héros (personnalités émergentes).
\item \jeu{Battle Brothers}, \jeu{Wildermyth} : la campagne \emph{est} l'histoire qu'on retient.
\end{itemize}
\end{optionbox}

\begin{optionbox}{Option C7 --- La tension sans RNG ni perte : l'information parfaite}
\jeu{Into the Breach} : les attaques ennemies sont \emph{télégraphiées}, le défi est un puzzle (protéger la grille en N tours). Tension maximale, zéro grind, zéro coup du sort. Une autre façon d'« avoir envie de bien jouer » sans punir par la perte.
\end{optionbox}

\begin{antipat}
\textbf{Tout réversible.} Save-scum, respec gratuite illimitée, undo systématique : la tension et le poids des choix s'effondrent. La réversibilité est une \emph{option de confort} (cf. C1 optionnelle, A8) à doser, pas un défaut par défaut.
\end{antipat}

% ======================================================================
\section{Synthèse : principes transférables}

\begin{enumerate}[leftmargin=1.4em]
\item \textbf{Pousser vers l'avant (une horloge) $>$ laisser farmer.} Une pression temporelle/systémique tue le grind à la source.
\item \textbf{Le choix avant l'accumulation.} Bifurcations, exclusivité, irréversibilité : chaque palier ferme des portes.
\item \textbf{La rareté crée la décision.} Récupération limitée, économie tendue, paie qui scale, coûts d'opportunité.
\item \textbf{L'enjeu crée le sens.} Perte permanente (ou séquelles), attrition, échec qui fait avancer.
\item \textbf{L'identité avant les stats.} Personnages nommés/customisés, historique émergent : on s'attache à des gens, pas à des chiffres.
\item \textbf{Curer, ne pas agrandir.} Cartes à nœuds, régions cloisonnées denses, golden path fait-main au milieu d'optionnel.
\item \textbf{La puissance = de nouveaux verbes, pas de plus gros nombres.} Horizontal $>$ vertical ; l'XP n'est qu'une option de montée parmi d'autres.
\item \textbf{Les couches se renforcent.} La même ressource qui sert l'enjeu (paie, Renown, Karma, slots de repos) doit aussi être le levier de décision : c'est l'imbrication macro/micro/tension qui rend une campagne « pleine ».
\end{enumerate}

\begin{principe}
\textbf{En une phrase.} La progression devient motivante quand \emph{monter, c'est décider} (choix exclusifs), \emph{continuer, c'est risquer} (attrition/perte) et \emph{explorer, c'est arbitrer} (carte décisionnelle curée) --- l'accumulation de stats n'étant que le décor de ces trois moteurs.
\end{principe}

\subsection*{Tableau d'aiguillage : quel mécanisme pour quel objectif}
\begin{center}\small
\begin{tabular}{p{5.3cm}p{9.2cm}}
\hline
\textbf{Si tu veux\dots} & \textbf{Regarde l'option} \\
\hline
Tuer le grind d'XP & A1 (horloge), B2 (pas d'XP), C3 (rareté/paie) \\
Tuer la linéarité & A2/A3 (nœuds/geoscape), C5 (choix exclusifs), B1/B6 (builds divergents) \\
Éviter le monde trop grand & A2 (nœuds), A5 (régions cloisonnées), A7 (scaling), densité $>$ taille \\
Donner du poids émotionnel & C1/C2 (perte/séquelles), C6 (récit émergent), B9 (identité) \\
Garder la boucle micro vivante & B1 (exclusivité), B7 (monnaie-dilemme), B8 (loot à malus) \\
Accessibilité / public large & A8 (pas de pression), C1 optionnelle, C7 (information parfaite) \\
\hline
\end{tabular}
\end{center}

% ======================================================================
\section*{Annexe --- sources \& fiabilité}
\addcontentsline{toc}{section}{Annexe --- sources \& fiabilité}
\small
\textbf{Sources de développeur (\dv) :} Darkest Dungeon, \emph{Affliction System Deep Dive} (Game Developer) ; XCOM, \emph{Why we get attached to our soldiers} + interview producteur (Game Developer / PC Gamer) ; Crusader Kings II, \emph{Emergent Stories} (GDC Vault, Fåhraeus) ; Invisible Inc., \emph{Alarm Systems Deep Dive} (Game Developer) ; Wildermyth (interviews / presse design). \\[3pt]
\textbf{Mécaniques confirmées (wikis / guides communautaires) :} XCOM EU/2/WOTC, Phoenix Point, Slay the Spire, FTL, Into the Breach, Hades, Battle Brothers, Wartales, Jagged Alliance 2/3, Bannerlord, Divinity: OS2, Pathfinder: WOTR/Kingmaker, Wasteland 2/3, Shadowrun Dragonfall/Hong Kong, The Banner Saga, Fire Emblem, King's Bounty. \\[3pt]
\textbf{Analyse de design (non sourcée) :} tous les encarts \emph{Compromis}, \emph{Principe}, \emph{Anti-pattern} et la synthèse transversale. \\[3pt]
\textit{Limite honnête :} une partie des extractions web a été contrainte par un proxy bloquant certaines pages (lecture sur extraits de recherche plutôt qu'intégrale) ; les faits structurels sont robustes, les chiffres exacts (seuils, pourcentages) peuvent varier selon patch/difficulté.

\end{document}
